\subsection{Semi-local rings}

PROPOSITION 16. Let A be a ring. The following properties are equivalent:

\begin{enumerate}
\def\labelenumi{(\alph{enumi})}
\item
  the set of maximal ideals of \(\mathbf{A}\) is finite;
\item
  the quotient of A by its Jacobson radical is the direct composition of a finite number of fields.
\end{enumerate}

Suppose that the quotient of A by its Jacobson radical \(\Re\) is a direct composition of a finite number of fields. Then A/ \(\Re\) possesses only a finite number of ideals and \(a\ fortiori\) only a finite number of maximal ideals. As every maximal ideal contains \(\Re\) (Algebra, Chapter VIII, § 6, no. 2, Definition 2), the maximal ideals of A are the inverse images of the maximal ideals of A/ \(\Re\) under the canonical homomorphism A \(\rightarrow\) A/ \(\Re\) ; hence they are finite in number.

Conversely, suppose that A has only a finite number of distinct maximal ideals \(m_{1}, \ldots, m_{n}\) . The \(A/m_{i}\) are fields and it follows from § 1, no. 2, Proposition 5 that the canonical mapping \(A \to \prod_{i=1}^{n} A/m_{i}\) is surjective; as its kernel \(\bigcap_{i=1}^{n} m_{i}\) is the Jacobson radical R (Algebra, Chapter VIII, § 6, no. 2, Definition 2), A/R is isomorphic to \(\prod_{i=1}^{n} A/m_{i}\) .

DEFINITION 4. A ring is called semi-local if it satisfies the equivalent conditions (a), (b) of Proposition 16.

Examples. Every local ring is semi-local. Every quotient of a semi-local ring is semi-local. Every finite product of semi-local rings is semi-local.

* If A is a Noetherian semi-local ring and B is an A-algebra which is a finitely generated A-module, then B is semi-local (Chapter IV, § 2, no. 5, Corollary 3 to Proposition 9). *

Another example, generalizing the construction of the local rings \(A_{p}\) , is provided by the following proposition:

PROPOSITION 17. Let A be a ring and \(p_1, \ldots, p_n\) prime ideals of A. We write \(S = \bigcap_{i=1}^{n} (A - p_i) = A - \bigcup_{i=1}^{n} p_i\) .

\begin{enumerate}
\def\labelenumi{(\alph{enumi})}
\item
  The ring \(\mathbf{S}^{-1}\mathbf{A}\) is semi-local; if \(\mathfrak{q}_1, \ldots, \mathfrak{q}_r\) are the distinct maximal elements (with respect to inclusion) of the set of \(\mathfrak{p}_i\) , the maximal ideals of \(\mathbf{S}^{-1}\mathbf{A}\) are the \(\mathbf{S}^{-1}\mathfrak{q}_j\) ( \(1 \leqslant j \leqslant r\) ) and these ideals are distinct.
\item
  The ring \(\mathbf{A}_{\mathfrak{p}_i}\) is canonically isomorphic to \((\mathbf{S}^{-1}\mathbf{A})_{\mathbf{S}^{-1}\mathfrak{p}_i}\) for \(1 \leqslant i \leqslant n\) .
\item
  If \(\mathbf{A}\) is an integral domain, then \(\mathbf{S}^{-1}\mathbf{A} = \bigcap_{i=1}^{n}\mathbf{A}_{p_i}\) in the field of fractions of \(\mathbf{A}\) .
\end{enumerate}

(a) The ideals of A not meeting S are the ideals contained in the union of the \(p_{i}\) and hence in at least one of the \(p_{i}\) (§ 1, no. 1, Proposition 2); the \(q_{j}\) are therefore the maximal elements of the set of ideals not meeting S; consequently, the \(S^{-1}q_{j}\) are the maximal ideals of \(S^{-1}A\) by § 2, no. 5, Proposition 11 (ii).

(b) is a special case of § 2, no. 5, Proposition 11 (iii).

(c) Suppose that A is an integral domain. If \(p_{i} \subset p_{k}\) , then \(A_{p_{i}} \supset A_{p_{k}}\) ; to prove (c), we may therefore suppose that no two \(p_{i}\) are comparable. Then it follows from (a) and no. 3, Corollary 4 to Theorem 1 that \(S^{-1}A = \bigcap_{i=1}^{n} (S^{-1}A)_{S^{-1}p_{i}}\) ; whence (c) by virtue of (b).

If A is an integral domain, so is \(S^{-1}A\) and Proposition 17 then provides an example of a semi-local ring which is not a direct composition of local rings (cf.~Chapter III, § 2, no. 13).

COROLLARY. Let A be an integral domain and \(p_{1}, \ldots, p_{n}\) prime ideals of A, no two of which are comparable with respect to inclusion. If \(A = \bigcap_{i=1}^{n} A_{p_{i}}\) in the field of fractions of A, the maximal ideals of A are \(p_{1}, \ldots, p_{n}\) .

Setting \(S = \bigcap_{i=1}^{n} (A - p_i), S^{-1} A = A\) by Proposition 17 (c); hence the elements of \(S\) are invertible in \(A\) and \(S^{-1} p_i = p_i\) for all \(i\) . Our assertion then follows by virtue of Proposition 17 (a).

